\documentclass[twocolumn,showpacs,preprintnumbers,amsmath,amssymb,prd]{revtex4-2}

\usepackage{graphicx}
\usepackage{amsmath}
\usepackage{amssymb}
\usepackage{amsthm}
\usepackage{hyperref}

% Theorem environments
\newtheorem{theorem}{Theorem}[section]
\newtheorem{lemma}[theorem]{Lemma}
\newtheorem{corollary}[theorem]{Corollary}
\newtheorem{definition}[theorem]{Definition}
\newtheorem{proposition}[theorem]{Proposition}
\newtheorem{remark}[theorem]{Remark}

\begin{document}

\title{Geometric Unification of Fundamental Constants:\\
Deriving $\alpha$, $\sin^2\theta_W$, and $a_0$ from Cube Topology}

\author{Carl Zimmerman}
\email{carl@zimmerman.com}
\affiliation{Independent Researcher}

\date{\today}

\begin{abstract}
We present a geometric framework for deriving fundamental physical constants from the topology of the cube. Starting from three axioms---that spacetime is fundamentally discrete, that the unit cell is the binary cube, and that physical theories must satisfy consistency conditions---we derive relationships for the fine structure constant ($\alpha^{-1} = 4Z^2 + 3 = 137.04$), the Weinberg angle ($\sin^2\theta_W = 3/13 = 0.231$), and the MOND acceleration scale ($a_0 = cH_0/Z$), where $Z = 2\sqrt{8\pi/3} = 5.79$ and $Z^2 = 32\pi/3$. The coefficient 4 equals the Gauss-Bonnet curvature factor for the cube, while the coefficient 3 arises from a lattice index condition requiring balance between fermionic and gauge degrees of freedom. These results suggest a geometric closure in which coupling constants are determined by topology rather than being free parameters.
\end{abstract}

\pacs{11.30.-j, 12.10.-g, 04.50.Kd, 02.40.-k}

\maketitle

%==============================================================================
\section{Introduction}
%==============================================================================

The Standard Model of particle physics contains approximately 19 free parameters, including coupling constants that must be measured experimentally. The question ``Why is $\alpha \approx 1/137$?'' has remained open since the fine structure constant was first measured. We present a geometric framework that constrains these constants.

The key hypothesis is that if spacetime is fundamentally discrete at the Planck scale with cubic unit cells, then consistency conditions on the cube lattice constrain the coupling constants to specific values.

\begin{remark}[Scope and Claims]
This paper presents a mathematical framework with suggestive numerical matches. We claim the framework is \emph{geometrically consistent} and \emph{numerically predictive}, but acknowledge that the physical mechanism connecting cube topology to gauge couplings requires further development.
\end{remark}

%==============================================================================
\section{Axioms and Framework}
%==============================================================================

\begin{definition}[The Fundamental Axioms]
We assume:
\begin{enumerate}
\item[\textbf{A1.}] \textbf{Discreteness}: Spacetime is fundamentally discrete at the Planck scale, with a regular lattice structure.
\item[\textbf{A2.}] \textbf{Binary Structure}: The fundamental unit cell has binary vertex coordinates, i.e., vertices at $(x,y,z)$ with $x,y,z \in \{0,1\}$.
\item[\textbf{A3.}] \textbf{Consistency}: Physical theories defined on this lattice must satisfy index-theoretic consistency conditions analogous to anomaly cancellation.
\end{enumerate}
\end{definition}

%==============================================================================
\section{The Cube Uniqueness Theorem}
%==============================================================================

\begin{theorem}[Cube Uniqueness]
\label{thm:cube}
The cube is the unique 3-dimensional convex polytope satisfying:
\begin{enumerate}
\item All vertices have binary coordinates $(x,y,z) \in \{0,1\}^3$
\item The polytope tiles $\mathbb{R}^3$ by translation
\item The number of vertices equals $2^{\dim}$
\end{enumerate}
\end{theorem}

\begin{proof}
\textbf{(1) Binary vertices:} The set $\{0,1\}^3$ contains exactly $2^3 = 8$ points. The convex hull of these 8 points is the unit cube $[0,1]^3$. This is verified by noting that each vertex is extremal (cannot be written as a convex combination of others).

\textbf{(2) Space-filling:} Among the five Platonic solids, only the cube tiles $\mathbb{R}^3$ by translation. For any convex polytope to tile space, copies meeting along faces must fill $360°$ around each edge. The dihedral angle of the cube is exactly $90°$, and $360°/90° = 4$ (an integer). For comparison:
\begin{itemize}
\item Tetrahedron: $70.5°$ ($360°/70.5° = 5.1$, not an integer)
\item Octahedron: $109.5°$ ($360°/109.5° = 3.3$, not an integer)
\end{itemize}

\textbf{(3) Dimensional matching:} The $n$-cube (hypercube in $n$ dimensions) has exactly $2^n$ vertices. For $n=3$: $2^3 = 8$. This is the defining property of the hypercube family, not shared by other regular polytopes.

The intersection of these three criteria uniquely selects the cube. \qed
\end{proof}

\begin{definition}[Cube Invariants]
The cube has the following topological invariants:
\begin{align}
V &= 8 \quad \text{(vertices)} \\
E &= 12 \quad \text{(edges)} \\
F &= 6 \quad \text{(faces)} \\
D &= 4 \quad \text{(space diagonals)} \\
\chi &= V - E + F = 2 \quad \text{(Euler characteristic)}
\end{align}
We define the symbolic names: $\text{CUBE} \equiv V = 8$, $\text{GAUGE} \equiv E = 12$, $\text{FACES} \equiv F = 6$, $\text{BEKENSTEIN} \equiv D = 4$.
\end{definition}

%==============================================================================
\section{The Lattice Index Theorem}
%==============================================================================

In lattice gauge theory, consistency requires balance between fermionic and bosonic degrees of freedom. We formalize this as follows.

\begin{theorem}[Lattice Index Condition]
\label{thm:index}
On a cubic lattice where fermions occupy vertices and gauge fields occupy edges, the index-theoretic consistency condition:
\begin{equation}
\text{(fermionic modes)} = \text{(gauge modes)}
\end{equation}
uniquely determines $N_{\text{gen}} = 3$ if we require:
\begin{equation}
V \times N_{\text{gen}} = E \times 2
\end{equation}
\end{theorem}

\begin{proof}
On the cube lattice:
\begin{itemize}
\item Fermion fields are defined at vertices, with $N_{\text{gen}}$ species per vertex
\item Total fermionic content: $V \times N_{\text{gen}} = 8 \times N_{\text{gen}}$
\item Gauge fields are defined on edges, each with 2 polarizations/orientations
\item Total gauge content: $E \times 2 = 12 \times 2 = 24$
\end{itemize}

The balance condition $8 N_{\text{gen}} = 24$ has unique solution:
\begin{equation}
\boxed{N_{\text{gen}} = 3}
\end{equation}
\qed
\end{proof}

\begin{remark}[Connection to Standard Anomaly Cancellation]
The condition $V \times N_{\text{gen}} = E \times 2$ is analogous to, but distinct from, the standard chiral anomaly cancellation $\text{Tr}(T_a\{T_b, T_c\}) = 0$. We interpret it as an index-theoretic constraint on the lattice, related to the Atiyah-Singer index theorem in the continuum limit. A rigorous derivation of this connection remains an open problem.
\end{remark}

\begin{corollary}[Emergent Time Dimension]
\label{cor:time}
Defining $N_{\text{time}} \equiv D - N_{\text{gen}} = 4 - 3 = 1$, we obtain exactly one time dimension, consistent with observation.
\end{corollary}

%==============================================================================
\section{The Geometric Scale $Z^2$}
%==============================================================================

\begin{definition}[The Geometric Scale]
We define the fundamental geometric scale:
\begin{equation}
Z^2 \equiv V \times \frac{4\pi}{3} = 8 \times \frac{4\pi}{3} = \frac{32\pi}{3} \approx 33.51
\end{equation}
where $4\pi/3$ is the volume of a unit sphere (radius 1).
\end{definition}

\begin{remark}
The quantity $4\pi/3$ is the natural 3-dimensional solid normalization factor, appearing in the volume formula $V = \frac{4}{3}\pi r^3$. It represents the ``spherical content'' of the cube's embedding space.
\end{remark}

\begin{lemma}
\label{lem:Z}
$Z = \sqrt{Z^2} = \sqrt{32\pi/3} = 2\sqrt{8\pi/3} \approx 5.7888$
\end{lemma}

%==============================================================================
\section{The Gauss-Bonnet Connection}
%==============================================================================

\begin{theorem}[Gauss-Bonnet for the Cube]
\label{thm:gaussbonnet}
The total Gaussian curvature of the cube surface satisfies:
\begin{equation}
\int_{\partial \text{cube}} K \, dA = 4\pi = D \times \pi
\end{equation}
where $D = 4$ is the number of space diagonals.
\end{theorem}

\begin{proof}
The Gauss-Bonnet theorem states that for any closed surface $S$:
\begin{equation}
\int_S K \, dA = 2\pi\chi(S)
\end{equation}

For the cube surface:
\begin{itemize}
\item Each face is flat: $K = 0$ on faces
\item Curvature is concentrated at vertices (conical singularities)
\item At each vertex, three faces meet at right angles
\item Angle sum at vertex: $3 \times \frac{\pi}{2} = \frac{3\pi}{2}$
\item Angle deficit: $\delta = 2\pi - \frac{3\pi}{2} = \frac{\pi}{2}$
\end{itemize}

Total curvature from 8 vertices:
\begin{equation}
\int K \, dA = 8 \times \frac{\pi}{2} = 4\pi
\end{equation}

Verification: $2\pi\chi = 2\pi \times 2 = 4\pi$ \checkmark

The ``Gauss-Bonnet factor'' is:
\begin{equation}
\frac{1}{\pi}\int K \, dA = 4 = D = \text{BEKENSTEIN}
\end{equation}
\qed
\end{proof}

\begin{remark}[Interpretive Note]
The equality between the Gauss-Bonnet factor (4) and the number of space diagonals (4) is a mathematical fact about the cube. We \emph{conjecture} that this equality is physically significant for determining the electromagnetic coupling, but this connection requires further theoretical development.
\end{remark}

%==============================================================================
\section{The Fine Structure Constant}
%==============================================================================

\begin{theorem}[The $\alpha$ Formula]
\label{thm:alpha}
The fine structure constant satisfies:
\begin{equation}
\boxed{\alpha^{-1} = D \times Z^2 + N_{\text{gen}} = 4Z^2 + 3 = \frac{128\pi + 9}{3}}
\end{equation}
with numerical value $137.0413...$, compared to the observed $137.0360...$
\end{theorem}

\begin{proof}
We decompose $\alpha^{-1}$ into geometric and quantum contributions:

\textbf{1. Geometric Contribution:}
The coefficient $D = 4$ is the Gauss-Bonnet factor (Theorem~\ref{thm:gaussbonnet}). The factor $Z^2 = 32\pi/3$ is the geometric scale (Definition 5.1). Their product:
\begin{equation}
\alpha^{-1}_{\text{geom}} = D \times Z^2 = 4 \times \frac{32\pi}{3} = \frac{128\pi}{3} \approx 134.04
\end{equation}

\textbf{2. Quantum Correction:}
The number of fermion generations $N_{\text{gen}} = 3$ (Theorem~\ref{thm:index}) contributes an additive correction:
\begin{equation}
\Delta\alpha^{-1} = N_{\text{gen}} = 3
\end{equation}

\textbf{3. Total:}
\begin{equation}
\alpha^{-1} = \frac{128\pi}{3} + 3 = \frac{128\pi + 9}{3} = 137.0413...
\end{equation}

\textbf{Comparison with observation:}
\begin{align}
\alpha^{-1}_{\text{predicted}} &= 137.0413 \nonumber\\
\alpha^{-1}_{\text{observed}} &= 137.0360 \nonumber\\
\text{Deviation} &= 0.0039\%
\end{align}
\qed
\end{proof}

\begin{remark}[Physical Interpretation]
The additive form $\alpha^{-1} = (\text{geometric}) + (\text{quantum})$ is reminiscent of RG running, where $\alpha^{-1}(\mu) = \alpha^{-1}(\mu_0) + \beta \ln(\mu/\mu_0)$. However, the exact coefficient +3 (rather than a logarithmic function) suggests this is a boundary/topological effect rather than standard RG running. This interpretation requires further investigation.
\end{remark}

\begin{corollary}[Equivalent Forms]
\begin{align}
\alpha^{-1} &= 4Z^2 + 3 \\
&= \frac{D^2 \times V \times \pi + N_{\text{gen}}^2}{N_{\text{gen}}} \\
&= \frac{16 \times 8 \times \pi + 9}{3} = \frac{128\pi + 9}{3}
\end{align}
\end{corollary}

%==============================================================================
\section{The Weinberg Angle}
%==============================================================================

\begin{theorem}[The Weinberg Angle Formula]
\label{thm:weinberg}
The Weinberg angle satisfies:
\begin{equation}
\boxed{\sin^2\theta_W = \frac{N_{\text{gen}}}{E + N_{\text{time}}} = \frac{3}{12 + 1} = \frac{3}{13}}
\end{equation}
with numerical value $0.2308$, compared to observed $0.2312$.
\end{theorem}

\begin{proof}
\textbf{Gauge Group Structure:}
The Standard Model gauge group SU(3)$_C \times$SU(2)$_L \times$U(1)$_Y$ has:
\begin{align}
\dim(\text{SU}(3)) &= 8 = V \quad \text{(cube vertices)} \nonumber\\
\dim(\text{SU}(2)) &= 3 = N_{\text{gen}} \nonumber\\
\dim(\text{U}(1)) &= 1 = N_{\text{time}} \nonumber\\
\text{Total} &= 12 = E \quad \text{(cube edges)}
\end{align}

\textbf{Electroweak Sector:}
The electroweak gauge group SU(2)$_L \times$U(1)$_Y$ has dimension $3 + 1 = 4 = D$.

After symmetry breaking, U(1)$_{\text{EM}}$ survives, and the total ``electroweak + electromagnetic'' structure has effective dimension:
\begin{equation}
E + N_{\text{time}} = 12 + 1 = 13
\end{equation}

\textbf{The Mixing Angle:}
The weak isospin contribution relative to the total:
\begin{equation}
\sin^2\theta_W = \frac{N_{\text{gen}}}{E + N_{\text{time}}} = \frac{3}{13} = 0.2308
\end{equation}

\textbf{Comparison:}
\begin{align}
\sin^2\theta_W^{\text{predicted}} &= 0.2308 \nonumber\\
\sin^2\theta_W^{\text{observed}} &= 0.2312 \nonumber\\
\text{Deviation} &= 0.19\%
\end{align}

\textbf{GUT Consistency:}
At the GUT scale, SU(5) unification predicts $\sin^2\theta_W = 3/8$. The ratio:
\begin{equation}
\frac{\sin^2\theta_W(\text{low})}{\sin^2\theta_W(\text{GUT})} = \frac{3/13}{3/8} = \frac{8}{13} \approx 0.615
\end{equation}
This is consistent with Standard Model RG evolution from $M_{\text{GUT}}$ to $M_Z$.
\qed
\end{proof}

%==============================================================================
\section{The MOND Acceleration Scale}
%==============================================================================

\begin{theorem}[MOND Scale Derivation]
\label{thm:mond}
The MOND acceleration scale satisfies:
\begin{equation}
\boxed{a_0 = \frac{cH_0}{Z} = \frac{c\sqrt{G\rho_c}}{2}}
\end{equation}
where $Z = 2\sqrt{8\pi/3} \approx 5.79$.
\end{theorem}

\begin{proof}
\textbf{Step 1: Critical Density.}
From the Friedmann equation:
\begin{equation}
H^2 = \frac{8\pi G}{3}\rho_c \quad \Rightarrow \quad \rho_c = \frac{3H^2}{8\pi G}
\end{equation}

\textbf{Step 2: Natural Acceleration Scale.}
Dimensional analysis gives $[c\sqrt{G\rho}] = $ acceleration:
\begin{equation}
a_{\text{nat}} = c\sqrt{G\rho_c} = c\sqrt{\frac{3H^2}{8\pi}} = \frac{cH}{\sqrt{8\pi/3}}
\end{equation}

\textbf{Step 3: Horizon Thermodynamics.}
The de Sitter horizon mass from Bekenstein-Hawking thermodynamics:
\begin{equation}
M_H = \frac{c^3}{2GH}
\end{equation}
The factor of 2 arises from the horizon entropy bound.

\textbf{Step 4: Combined Scale.}
\begin{equation}
a_0 = \frac{a_{\text{nat}}}{2} = \frac{cH}{2\sqrt{8\pi/3}} = \frac{cH}{Z}
\end{equation}

\textbf{Numerical Evaluation} ($H_0 = 67.4$ km/s/Mpc):
\begin{equation}
a_0 = 1.13 \times 10^{-10} \text{ m/s}^2
\end{equation}
Observed: $a_0 \approx (1.2 \pm 0.1) \times 10^{-10}$ m/s$^2$.

Agreement within $\sim 6\%$, comparable to $H_0$ measurement uncertainty.
\qed
\end{proof}

%==============================================================================
\section{The Geometric Closure Theorem}
%==============================================================================

\begin{theorem}[Geometric Closure]
\label{thm:closure}
All coefficients appearing in the formulas for $\alpha^{-1}$, $\sin^2\theta_W$, and $a_0$ are determined by cube topology:
\begin{center}
\begin{tabular}{|c|l|l|}
\hline
Value & Symbol & Source \\
\hline
8 & $V$ (CUBE) & Cube vertices \\
12 & $E$ (GAUGE) & Cube edges \\
6 & $F$ (FACES) & Cube faces \\
4 & $D$ (BEKENSTEIN) & Space diagonals = Gauss-Bonnet factor \\
3 & $N_{\text{gen}}$ & Index condition: $8 \times 3 = 24 = 12 \times 2$ \\
1 & $N_{\text{time}}$ & $D - N_{\text{gen}} = 4 - 3$ \\
$32\pi/3$ & $Z^2$ & $V \times (4\pi/3)$ \\
\hline
\end{tabular}
\end{center}
There are no free parameters except measured cosmological inputs ($c$, $H_0$).
\end{theorem}

%==============================================================================
\section{Discussion}
%==============================================================================

\subsection{Summary of Results}

\begin{table}[h]
\centering
\begin{tabular}{|l|c|c|c|}
\hline
Quantity & Formula & Predicted & Observed \\
\hline
$\alpha^{-1}$ & $4Z^2 + 3$ & 137.041 & 137.036 \\
$\sin^2\theta_W$ & $3/13$ & 0.2308 & 0.2312 \\
$a_0$ & $cH_0/Z$ & $1.13 \times 10^{-10}$ & $1.2 \times 10^{-10}$ \\
$N_{\text{gen}}$ & $24/8$ & 3 & 3 \\
\hline
\end{tabular}
\caption{Comparison of predictions with observations.}
\end{table}

\subsection{Strengths of the Framework}

\begin{enumerate}
\item \textbf{Numerical accuracy}: All predictions match observations to better than 1\% (except $a_0$, which matches within $H_0$ uncertainty).
\item \textbf{No fitting}: The formulas contain no adjustable parameters.
\item \textbf{Geometric origin}: All coefficients derive from cube topology.
\item \textbf{Internal consistency}: The framework predicts $N_{\text{gen}} = 3$ and $N_{\text{time}} = 1$ correctly.
\end{enumerate}

\subsection{Open Questions and Limitations}

\begin{enumerate}
\item \textbf{Physical mechanism}: We have not derived \emph{why} gauge couplings should depend on cube topology. The connection remains conjectural.
\item \textbf{Index condition}: The lattice balance $V \times N_{\text{gen}} = E \times 2$ resembles but is not identical to standard anomaly cancellation. A rigorous derivation from index theory is needed.
\item \textbf{Additive structure}: Why $\alpha^{-1} = 4Z^2 + 3$ rather than multiplicative? The additive form suggests boundary effects but lacks derivation.
\item \textbf{Weinberg angle +1}: The ``+1'' in the denominator (giving 13 instead of 12) needs stronger justification.
\end{enumerate}

\subsection{Testable Predictions}

\begin{enumerate}
\item \textbf{MOND evolution}: If $a_0 = cH_0/Z$, then $a_0(z) \propto H(z)$, giving:
\begin{equation}
a_0(z) = a_0(0) \times \sqrt{\Omega_m(1+z)^3 + \Omega_\Lambda}
\end{equation}
This is testable with high-redshift kinematic observations.
\item \textbf{No fourth generation}: The framework predicts $N_{\text{gen}} = 3$ exactly.
\item \textbf{Precision constants}: Any deviation from $\alpha^{-1} = 4Z^2 + 3$ or $\sin^2\theta_W = 3/13$ would falsify the framework.
\end{enumerate}

%==============================================================================
\section{Conclusion}
%==============================================================================

We have presented a geometric framework relating fundamental constants to cube topology:
\begin{align}
\alpha^{-1} &= 4Z^2 + 3 = 137.04 \nonumber\\
\sin^2\theta_W &= 3/13 = 0.231 \nonumber\\
a_0 &= cH_0/Z
\end{align}

The framework achieves remarkable numerical accuracy without free parameters. Every coefficient is either a topological invariant of the cube (4, 8, 12, 3, 1) or derived from consistency conditions (lattice index balance).

While the physical mechanism connecting cube topology to gauge couplings remains to be established, the numerical success suggests this direction warrants further investigation.

\begin{acknowledgments}
The author thanks colleagues for critical discussions on mathematical rigor.
\end{acknowledgments}

\begin{thebibliography}{99}

\bibitem{Milgrom1983} M. Milgrom, Astrophys. J. \textbf{270}, 365 (1983).

\bibitem{Milgrom2020} M. Milgrom, arXiv:2001.09729 (2020).

\bibitem{Famaey2012} B. Famaey and S. McGaugh, Living Rev. Rel. \textbf{15}, 10 (2012).

\bibitem{Planck2020} Planck Collaboration, Astron. Astrophys. \textbf{641}, A6 (2020).

\bibitem{McGaugh2016} S. McGaugh et al., Phys. Rev. Lett. \textbf{117}, 201101 (2016).

\bibitem{Verlinde2017} E. Verlinde, SciPost Phys. \textbf{2}, 016 (2017).

\bibitem{Chern1944} S.-S. Chern, Ann. Math. \textbf{45}, 747 (1944).

\bibitem{AtiyahSinger1968} M. Atiyah and I. Singer, Ann. Math. \textbf{87}, 484 (1968).

\bibitem{PDG2022} R.L. Workman et al. (Particle Data Group), Prog. Theor. Exp. Phys. 2022, 083C01 (2022).

\end{thebibliography}

\end{document}
